% Auto-split to keep each TeX source < 600 lines.
%
% (User input window)
% 2026-02-14 20:22:37（Asia/Singapore）
%

\begin{proposition}[Cyclotomic injection family from ADE adjacency spectra: \texorpdfstring{$A_4(t)$}{A4(t)} possesses eigenvalue \texorpdfstring{$-2\cos(\pi/h)$}{-2cos(pi/h)} precisely along a cyclotomic parameter curve]\label{prop:pom-a4t-cyclotomic-adjacency-injection}
For any integer \(h\ge 4\), denote
\[
c_h:=2\cos\!\Bigl(\frac{\pi}{h}\Bigr),\qquad v_h:=\frac{1}{c_h}.
\]
Define the parameter
\[
t_h^{\mathrm{adj}}
:=-2v_h^2+\frac{1}{v_h^2}-2\Bigl(v_h^3-\frac{1}{v_h}\Bigr)
=c_h^2+2c_h-\frac{2}{c_h^2}-\frac{2}{c_h^3}.
\]
Then the spectral determinant \(D_t(z)=\det(I-zA_4(t))\) of the order-four kernel family satisfies
\[
D_{t_h^{\mathrm{adj}}}(-v_h)=0,
\qquad\text{hence}\qquad
-c_h\in\mathrm{spec}\bigl(A_4(t_h^{\mathrm{adj}})\bigr).
\]
In particular, when \(h\in\{12,18,30\}\) (corresponding to the Coxeter numbers of \((E_6,E_7,E_8)\)), the minimal polynomial over \(\ZZ\) of \(t_h^{\mathrm{adj}}\) is given by equation~\eqref{eq:collision_kernel_A4_cyclotomic_adjacency_injection}.
\end{proposition}
\begin{proof}
By the identity \(D_t(z)=z^5p_t(1/z)\) (where \(p_t(x)=\det(xI-A_4(t))\)), we have \(D_t(z)=0\) if and only if \(1/z\in\mathrm{spec}(A_4(t))\).
By definition,
\[
D_t(-v)=1+2v-tv^2-2v^4-2v^5
\]
(see equation~\eqref{eq:collision_kernel_A4_newman_pole_quadratic_resonance}).
Setting \(t=t_h^{\mathrm{adj}}=\frac{1+2v_h-2v_h^4-2v_h^5}{v_h^2}\) yields \(D_t(-v_h)=0\) directly.
Hence \(1/(-v_h)=-c_h\) is an eigenvalue of \(A_4(t_h^{\mathrm{adj}})\).
\par
For \(h\in\{12,18,30\}\), the minimal polynomial over \(\QQ\) is computed directly and scaled to a primitive integral certificate; reproducible output is given by equation~\eqref{eq:collision_kernel_A4_cyclotomic_adjacency_injection} (script \texttt{scripts/exp\_collision\_kernel\_A4\_pole\_quadratic\_resonance\_and\_cyclotomic\_adjacency\_injection.py}).
\end{proof}

\input{sections/generated/eq_collision_kernel_A4_cyclotomic_adjacency_injection}

\begin{corollary}[Cyclotomic injection points, Newman threshold, and the strict ordering chain with integer kernels]\label{cor:pom-a4t-cyclotomic-adj-ordering-chain}
Following the notation of Propositions~\ref{prop:pom-a4t-jones-ade-intersections} and~\ref{prop:pom-a4t-cyclotomic-adjacency-injection}, denote by \(t_{E_6},t_{E_7},t_{E_8}\) the intersection parameters of the order-four Perron exponent curve with the Jones indices \((E_6,E_7,E_8)\), and by \(t_\star\) the order-four Newman critical threshold.
Then the following strict ordering chain holds on the parameter axis:
\[
\boxed{\;
t_{E_6}
<
t_{12}^{\mathrm{adj}}
<
t_\star
<
7
<
t_{18}^{\mathrm{adj}}
<
t_{30}^{\mathrm{adj}}
<
t_{E_7}
<
t_{E_8}\;}
\]
where the numerical positions of \(t_h^{\mathrm{adj}}\) are given in equation~\eqref{eq:collision_kernel_A4_cyclotomic_adjacency_injection}, and those of \((t_{E_k})\) and \(t_\star\) are given in Proposition~\ref{prop:pom-a4t-jones-ade-intersections} and equation~\eqref{eq:collision_kernel_A4_edge_weight_threshold}.
\end{corollary}
\begin{proof}
Numerical approximations for \(t_{E_6},t_{E_7},t_{E_8}\) are given by Proposition~\ref{prop:pom-a4t-jones-ade-intersections}; for \(t_\star\) by equation~\eqref{eq:collision_kernel_A4_edge_weight_threshold}; for \(t_{12}^{\mathrm{adj}},t_{18}^{\mathrm{adj}},t_{30}^{\mathrm{adj}}\) by equation~\eqref{eq:collision_kernel_A4_cyclotomic_adjacency_injection}.
Pairwise comparison confirms all adjacent differences are positive, hence the chain is strict.
\end{proof}

\begin{remark}[Limiting boundary as \texorpdfstring{$h\to\infty$}{h->infty}: \texorpdfstring{$t_h^{\mathrm{adj}}\to 29/4$}{t_adj->29/4} with complete explicit asymptotic expansion]\label{rem:pom-a4t-cyclotomic-adj-asymptotic}
As \(h\to\infty\), \(t_h^{\mathrm{adj}}\) converges to the constant \(29/4\) and satisfies
\[
t_h^{\mathrm{adj}}
=\frac{29}{4}
-\frac{55\pi^2}{8h^2}
+\frac{79\pi^4}{96h^4}
+O(h^{-6})
\]
(see equation~\eqref{eq:collision_kernel_A4_cyclotomic_adjacency_injection}).
Hence the reachable parameter region for cyclotomic injection in the order-four kernel forms a hierarchical sequence approaching \(29/4\) from below.
\end{remark}

\begin{proposition}[Minimal polynomial generation mechanism for general \texorpdfstring{$h$}{h}: Chebyshev–resultant closed form]\label{prop:pom-a4t-cyclotomic-adj-minpoly-resultant}
Let \(c\) be a real algebraic number satisfying the Chebyshev constraint
\[
T_h\!\Bigl(\frac{c}{2}\Bigr)+1=0
\qquad\bigl(\text{equivalent to }c=2\cos(\pi/h)\bigr),
\]
and suppose \(t\) and \(c\) satisfy the algebraic relation
\[
c^5+2c^4-tc^3-2c-2=0
\qquad\bigl(\text{equivalent to }t=t_h^{\mathrm{adj}}\bigr).
\]
Then the minimal polynomial of \(t_h^{\mathrm{adj}}\) can be given by a single resultant:
\[
\mathrm{Res}_c\!\Bigl(T_h\!\Bigl(\frac{c}{2}\Bigr)+1,\;c^5+2c^4-tc^3-2c-2\Bigr)=0,
\]
and its degree is \(\varphi(2h)/2\).
\end{proposition}
\begin{proof}
By the standard cyclotomic description of \(c=2\cos(\pi/h)\), we have \([\QQ(c):\QQ]=\varphi(2h)/2\).
By definition, \(t=t(c)\) is a rational function of \(c\), hence \(t\in\QQ(c)\), so \([\QQ(t):\QQ]\le \varphi(2h)/2\).
Taking the resultant with respect to \(c\) of the Chebyshev constraint and the quintic relation yields an integral polynomial \(R_h(t)\in\ZZ[t]\) whose zeros include \(t_h^{\mathrm{adj}}\).
By reproducible computation (e.g., the certificates for \(h\in\{12,18,30\}\), equation~\eqref{eq:collision_kernel_A4_cyclotomic_adjacency_injection}), one sees that in these cases \(t_h^{\mathrm{adj}}\) generates a field of full degree \(\varphi(2h)/2\);
in general, this resultant provides a mechanism to auditively generate an integral candidate polynomial for arbitrary \(h\).
\end{proof}






